% TOP_PROBLEM: {"id": 20000880, "title": "Restricted Roth progressions and their local obstructions", "category_id": 2, "category": {"id": 2, "name": "combinatorics", "display_name": "Combinatorics", "description": "Counting problems, graph theory, discrete structures.", "slug": "combinatorics", "order_index": 2, "created_at": "2026-07-31T15:26:25.671Z"}, "status": "partially_solved", "problem_number": "AIM-COMBINATORICS-0005", "set_id": 14}
% FIRST_POSTED: 2026-10-01
% ENTRY_KIND: statement_only
% UnsolvedMath ID 20000880; source catalogue code AIM-COMBINATORICS-0005
% !TeX program = lualatex
% Definition and sources only; this file is not an LLM solution attempt.
\documentclass[11pt,a4paper]{article}
\usepackage[margin=25mm]{geometry}
\usepackage{fontspec}
\setmainfont{lmroman10-regular.otf}[BoldFont=lmroman10-bold.otf,
 ItalicFont=lmroman10-italic.otf,BoldItalicFont=lmroman10-bolditalic.otf]
\usepackage{amsmath,amssymb,mathtools}
\usepackage{unicode-math}
\setmathfont{latinmodern-math.otf}
\AtBeginDocument{\let\setminus\smallsetminus}
\usepackage{xurl,hyperref}
\hypersetup{colorlinks=true,urlcolor=blue}
\setcounter{secnumdepth}{0}
\setlength{\parindent}{0pt}
\setlength{\parskip}{5pt}
\setlength{\emergencystretch}{3em}
\begin{document}
\section*{Restricted Roth progressions and their local obstructions}\label{20000880}
{\small UnsolvedMath ID 20000880 · AIM-COMBINATORICS-0005 · Combinatorics · AIM Workshop Problem Lists}
\subsection{Definitions and mathematical statement}
Write $[N]=\{1,\ldots,N\}$, and let $\log$ denote the natural logarithm.
A nonconstant three-term arithmetic progression in $[N]$ is a triple
$x,x+d,x+2d$ with integers $x,d$, $d\geq1$, and $1\leq x<x+2d\leq N$.
Here the permitted differences are the shifted primes
$\mathcal D=\{p-1:p\text{ is prime}\}$; in particular, the prime $2$ gives the permitted difference $1$.
All three terms must belong to the same set, and they are distinct.

The conjecture asks whether, for every real $C>0$, there exists an integer $N_0(C)$ such that for every $N\geq N_0(C)$ and every $A\subseteq[N]$,
\[
 |A|\geq\frac{N}{(\log N)^C}
 \quad\Longrightarrow\quad
 \exists\,p\text{ prime},\ x\in\mathbb Z:
 \{x,x+p-1,x+2(p-1)\}\subseteq A.
\]
Equivalently, if $r_{p-1}(N)$ denotes the largest cardinality of a subset of $[N]$ avoiding all these triples, prove
$r_{p-1}(N)<N/(\log N)^C$ for every fixed $C>0$ once $N$ is sufficiently large.
The threshold may depend on $C$, but the conclusion must hold uniformly over the sets $A$.
The restriction is on the common difference, not on the primality of the three terms themselves.
This is the prime-minus-one question from the workshop; variations involving sums of two squares are separate questions and do not alter the specified difference set.
\subsection{Short English statement}
Does every set denser than any fixed inverse power of the logarithm contain a three-term progression whose difference is one less than a prime?
\subsection{Sources}
\begin{itemize}
\item[S1] ULAM.AI and the UnsolvedMath Contributors, supplied problems.json, record 20000880 (AIM-COMBINATORICS-0005), AIM Workshop Problem Lists. Catalogue identity and source transcription; CC BY 4.0.
\url{https://huggingface.co/datasets/ulamai/UnsolvedMath}.
\item[S2] American Institute of Mathematics, High-dimensional phenomena in discrete analysis, item 2.2. Original workshop question underlying AIM-COMBINATORICS-0005.
\url{http://aimpl.org/highdimdiscrete/2/}.
\item[S3] Ben Green, Open problems, problem 11(ii). Quantitative shifted-prime progression context.
\url{https://people.maths.ox.ac.uk/greenbj/papers/open-problems.pdf}.
\end{itemize}
\end{document}
